% Options for packages loaded elsewhere
\PassOptionsToPackage{unicode}{hyperref}
\PassOptionsToPackage{hyphens}{url}
\documentclass[
  ignorenonframetext,
]{beamer}
\newif\ifbibliography
\usepackage{pgfpages}
\setbeamertemplate{caption}[numbered]
\setbeamertemplate{caption label separator}{: }
\setbeamercolor{caption name}{fg=normal text.fg}
\beamertemplatenavigationsymbolsempty
% remove section numbering
\setbeamertemplate{part page}{
  \centering
  \begin{beamercolorbox}[sep=16pt,center]{part title}
    \usebeamerfont{part title}\insertpart\par
  \end{beamercolorbox}
}
\setbeamertemplate{section page}{
  \centering
  \begin{beamercolorbox}[sep=12pt,center]{section title}
    \usebeamerfont{section title}\insertsection\par
  \end{beamercolorbox}
}
\setbeamertemplate{subsection page}{
  \centering
  \begin{beamercolorbox}[sep=8pt,center]{subsection title}
    \usebeamerfont{subsection title}\insertsubsection\par
  \end{beamercolorbox}
}
% Prevent slide breaks in the middle of a paragraph
\widowpenalties 1 10000
\raggedbottom
\AtBeginPart{
  \frame{\partpage}
}
\AtBeginSection{
  \ifbibliography
  \else
    \frame{\sectionpage}
  \fi
}
\AtBeginSubsection{
  \frame{\subsectionpage}
}
\usepackage{iftex}
\ifPDFTeX
  \usepackage[T1]{fontenc}
  \usepackage[utf8]{inputenc}
  \usepackage{textcomp} % provide euro and other symbols
\else % if luatex or xetex
  \usepackage{unicode-math} % this also loads fontspec
  \defaultfontfeatures{Scale=MatchLowercase}
  \defaultfontfeatures[\rmfamily]{Ligatures=TeX,Scale=1}
\fi
\usepackage{lmodern}
\ifPDFTeX\else
  % xetex/luatex font selection
\fi
% Use upquote if available, for straight quotes in verbatim environments
\IfFileExists{upquote.sty}{\usepackage{upquote}}{}
\IfFileExists{microtype.sty}{% use microtype if available
  \usepackage[]{microtype}
  \UseMicrotypeSet[protrusion]{basicmath} % disable protrusion for tt fonts
}{}
\makeatletter
\@ifundefined{KOMAClassName}{% if non-KOMA class
  \IfFileExists{parskip.sty}{%
    \usepackage{parskip}
  }{% else
    \setlength{\parindent}{0pt}
    \setlength{\parskip}{6pt plus 2pt minus 1pt}}
}{% if KOMA class
  \KOMAoptions{parskip=half}}
\makeatother
\setlength{\emergencystretch}{3em} % prevent overfull lines
\providecommand{\tightlist}{%
  \setlength{\itemsep}{0pt}\setlength{\parskip}{0pt}}
\usepackage{bookmark}
\IfFileExists{xurl.sty}{\usepackage{xurl}}{} % add URL line breaks if available
\urlstyle{same}
\hypersetup{
  hidelinks,
  pdfcreator={LaTeX via pandoc}}

\author{\texorpdfstring{}{}}
\date{}

\begin{document}

\begin{frame}{Abstract}
\protect\phantomsection\label{abstract}
No prior automated system tracks hypothesis-level evidence across the
full Active Inference and Free Energy Principle (FEP) literature. Manual
synthesis cannot keep pace with a field that has grown at a compound
annual rate of 16.99\% across 2005--2026, and the FEP's theoretical
generality has invited falsifiability critiques \citep{colombo2021free}
that only hypothesis-specific evidence profiling can address. Building
on the systematic literature analysis of Knight, Cordes, and Friedman
\citep{knight2022fep}---which pioneered manual annotation paired with
ontology-based analysis at the scale of hundreds of papers---we present
a computational meta-analysis framework that automates and scales this
approach.

The pipeline retrieves literature from arXiv, Semantic Scholar, and
OpenAlex, deduplicating \(N = 849\) papers via a canonical identifier
hierarchy (DOI \(>\) arXiv ID \(>\) Semantic Scholar ID \(>\) OpenAlex
ID). It classifies papers into a three-tier taxonomy spanning eight
categories: A (Core Theory), B (Tools \& Translation), and C
(Application Domains). An LLM-powered extraction system then evaluates
each abstract against eight core hypotheses, producing structured
nanopublications---each encoding directionality, a confidence score, and
natural-language reasoning---that populate an RDF-compatible knowledge
graph scored by a citation-weighted evidence function.

The resulting evidence landscape reveals a field where core theory
(Domain A) still dominates but tools development (Domain B)---including
pymdp, RxInfer.jl, and interpretable alternatives such as Free Energy
Projective Simulation---and application domains (Domain C) are
diversifying rapidly. Non-negative matrix factorization identifies five
latent topics that cross-cut the keyword taxonomy, and citation network
analysis exposes a sparse yet structured graph (1,678 intra-corpus
edges, 5.5\% reference resolution) anchored by pronounced hub papers.
Hypothesis scores range from strong consensus (H4 Predictive Coding, H5
Scalability) through moderate emerging support (H6 Clinical Utility, H7
Morphogenesis, H8 Language) to persistent contestation (H3 Markov
Blanket Realism)---though absolute score magnitudes are inflated by
publication bias and linguistic asymmetry in academic writing, making
relative rankings and temporal trajectories more reliable than point
estimates. By demonstrating that automated LLM-driven assertion
extraction can generate scalable, queryable representations of
scientific evidence, this work provides a reusable architecture for
\emph{living literature reviews}---continuously updated knowledge graphs
that track hypothesis-level consensus across rapidly evolving fields.

\textbf{Keywords:} Active Inference, Free Energy Principle,
meta-analysis, knowledge graph, nanopublications, bibliometrics,
hypothesis scoring, LLM extraction, computational neuroscience
\end{frame}

\end{document}
